
\section{Pinned Contract Evaluation}
\label{sec:rector_eval}



\subsection{Objective and pinned instrument}
\label{sec:pce_objective_instrument}

PCE measures modernization as \emph{progress against a fixed upgrade policy}, rather than as full functional correctness.
We encode the policy as a \emph{pinned} Rector configuration with a fixed Rector version, ruleset, and parameters.\cite{rector}
We treat this configuration as a migration contract.

We apply the contract twice.
First, we run it on the original benchmark.
For each file, the contract emits a set of upgrade \emph{rule types} that trigger, and this set defines that file's \emph{obligations} under the policy.
Second, we run the same contract on migrated outputs and measure (i) which baseline obligations no longer trigger (discharged obligations) and (ii) which new rule types now trigger (introduced obligations, meaning contract-visible regressions).
If Rector cannot analyze a migrated file (e.g., parser/runtime failure during dry-run), we exclude that file from PCE aggregates for that model and report analyzed coverage.

To focus on language migration, we score only \texttt{Rector\textbackslash Php*} rules.
We report additional diagnostics to contextualize contract outcomes (Section~\ref{sec:pce_validity_diagnostics}) and triangulate with PHPCompatibility (Section~\ref{sec:pce_phpcompat}).\cite{phpcompat}

\subsection{Signals: compliance breadth and residual-work proxy}
\label{sec:pce_primary_scores}

PCE reports two complementary signals.

\textbf{(1) Obligation discharge (primary).}
This measures breadth of compliance, meaning how much of the contract-defined obligation set is cleared.
Our unit of analysis is a \emph{file-level rule-type incidence}: a rule type counts at most once per file, regardless of how many times it occurs within that file.

\textbf{(2) Patch volume closure (secondary).}
This summarizes how much Rector dry-run patch volume remains after migration.
We use it only as a diff-magnitude proxy, not as developer effort.

% =========================
% Per-file metrics table
% (ICSE-style: compact, declarative, minimal prose)
% =========================
\begin{table}[t]
\centering
\caption{Per-file metrics recorded under the pinned contract.}
\label{tab:per_file_eval_metrics}
\small
\setlength{\tabcolsep}{4pt}
\renewcommand{\arraystretch}{1.12}
\begin{tabularx}{\columnwidth}{@{}>{\raggedright\arraybackslash}p{0.34\columnwidth} >{\raggedright\arraybackslash}X@{}}
\toprule
\textbf{Metric} & \textbf{Definition / description} \\
\midrule
Original triggers ($\mathcal{O}_i$, $O_i$) &
Rule \emph{types} triggered in file $i$ pre-migration; $O_i=|\mathcal{O}_i|$. \\
Remaining triggers ($\mathcal{R}_i$, $R_i$) &
Rule \emph{types} triggered post-migration; $R_i=|\mathcal{R}_i|$. \\
Removed / introduced ($S_i$, $I_i$) &
$S_i=|\mathcal{O}_i\setminus\mathcal{R}_i|$;\;
$I_i=|\mathcal{R}_i\setminus\mathcal{O}_i|$. \\
Net change ($\Delta_i$) &
$\Delta_i=S_i-I_i$ (negative: more types introduced). \\
Discharge rate ($\rho_i$) &
$\rho_i=\frac{S_i}{O_i}\times 100$. \\
Patch volume ($D_i$) &
Changed-line count in Rector dry-run patch (\texttt{diff\_line\_count}); residual edit-mass proxy. \\
Patch closure (\%) &
Closure relative to $D_i^{\mathrm{orig}}$ (WDC, MFDC). \\
Metadata &
LOC, size bin, identifiers; residual set $\mathcal{R}_i$ and milestone/rule breakdowns. \\
\bottomrule
\end{tabularx}
\end{table}
% =========================
% Outcome matrices table (syntax + loadability)
% =========================
\begin{table}[t]
\centering
\caption[Outcome matrices for syntax and loadability]{Outcome matrices for syntax lint and the isolated \emph{loadability tripwire} (regression-oriented: Orig.\ Pass $\rightarrow$ Migr.\ Fail).}
\label{tab:outcome_matrices}
\scriptsize
\setlength{\tabcolsep}{3pt}
\renewcommand{\arraystretch}{1.05}
\begin{subtable}[t]{0.46\columnwidth}
\centering
\caption{Syntax}
\label{tab:syntax_validation_outcomes}
\begin{tabular}{@{}lcc@{}}
\toprule
Outcome & Orig. & Migr. \\
\midrule
Both valid  & Pass & Pass \\
Fixed error & Fail & Pass \\
New error   & Pass & Fail \\
Both error  & Fail & Fail \\
\bottomrule
\end{tabular}
\end{subtable}
\hspace{6pt}
\begin{subtable}[t]{0.46\columnwidth}
\centering
\caption{Loadability}
\label{tab:runtime_validation_outcomes}
\begin{tabular}{@{}lcc@{}}
\toprule
Outcome & Orig. & Migr. \\
\midrule
Both pass     & Pass & Pass \\
Fixed failure & Fail & Pass \\
New failure   & Pass & Fail \\
Both fail     & Fail & Fail \\
\bottomrule
\end{tabular}
\end{subtable}
\end{table}

% =========================
% Structural diagnostics table
% =========================
\begin{table}[t]
\centering
\caption{Structural preservation signals.}
\label{tab:structural_completeness_signals}
\small
\setlength{\tabcolsep}{4pt}
\renewcommand{\arraystretch}{1.10}
\begin{tabularx}{\columnwidth}{@{}l X@{}}
\toprule
\textbf{Signal} & \textbf{Purpose} \\
\midrule
Token dist.\ sim. & Token-type frequency similarity (non-trivia). \\
Token seq.\ sim.  & Order-sensitive similarity on normalized token stream. \\
Decl.\ overlap    & Set overlap of namespaces/classes/functions/methods. \\
Line preservation & Code-line retention (excluding blank/comment lines). \\
Body preservation & Body token retention for matched declarations. \\
Signature stability & Parameter/modifier consistency for matched declarations. \\
Empty-body rate   & Fraction of matched bodies that are near-empty. \\
Trivial-stub rate & Fraction of matched bodies with shallow stub patterns. \\
\midrule
\multicolumn{2}{@{}p{\columnwidth}@{}}{\footnotesize\emph{Note:} Body/stub signals apply only when bodies exist; otherwise N/A.} \\
\bottomrule
\end{tabularx}
\end{table}

\subsubsection{Breadth: obligation discharge}
\label{sec:pce_breadth}

For file $i$, let $\mathcal{O}_i$ be the set of rule types that trigger on the original file under the pinned contract, and let $\mathcal{R}_i$ be the set that trigger on the migrated output.
Let $O_i = |\mathcal{O}_i|$ and $R_i = |\mathcal{R}_i|$.

We decompose change into three quantities:
\[
\begin{aligned}
S_i &= |\mathcal{O}_i \setminus \mathcal{R}_i|
\quad \text{(discharged types)}, \\
I_i &= |\mathcal{R}_i \setminus \mathcal{O}_i|
\quad \text{(introduced types)}, \\
\Delta_i &= S_i - I_i
\quad \text{(net change)}.
\end{aligned}
\]
Intuitively, $S_i$ answers: ``How many baseline rule types stopped triggering in this file?''
$I_i$ answers: ``How many new rule types did the migration introduce under the same policy?''

We define the per-file discharge rate as:
\[
\rho_i = \frac{S_i}{O_i}\times 100.
\]
All benchmark files satisfy $O_i>0$ by construction.
We report $I_i$ explicitly to surface contract-visible regressions.
These are not behavioral claims, only new triggers under the pinned policy.

\subsubsection{Secondary proxy: patch volume closure}
\label{sec:pce_depth}

Let $D_i$ be Rector's dry-run changed-line count for file $i$ on the migrated output, and let $D_i^{\mathrm{orig}}$ be the same quantity on the original file under the same pinned contract.
We summarize reduction in patch volume in two ways: (i) workload-weighted closure (WDC), which emphasizes high-volume files, and (ii) mean file-level closure (MFDC), which treats files uniformly.
If either quantity is missing or unreadable for a file (e.g., failed Rector analysis output), that file is excluded from numeric patch-volume aggregates.

\providecommand{\WDC}{\textnormal{WDC}}
\providecommand{\MFDC}{\mathrm{MFDC}}
{\footnotesize
\begin{align}
\WDC(m)
&= \left(1-\frac{\sum_i D_i^{m}}{\sum_i D_i^{\mathrm{orig}}}\right)\cdot 100,\\
\MFDC(m)
&= \frac{1}{N}\sum_i c_i,\\[-2pt]
c_i
&=
\begin{aligned}[t]
\left(1-\frac{D_i^{m}}{D_i^{\mathrm{orig}}}\right)\cdot 100 \;& (D_i^{\mathrm{orig}}>0)\\
100 \;& (D_i^{\mathrm{orig}}=0,\, D_i^{m}=0)\\
0 \;& (D_i^{\mathrm{orig}}=0,\, D_i^{m}>0)\;.
\end{aligned}
\end{align}
}
Patch volume closure can be inflated by stylistic rewrites or diff-shaping changes, so we use it only as supporting context for residual work, not as a correctness or effort estimate.



\subsection{Reporting and interpretation}
\label{sec:pce_interpretation}

For each file we record the obligation counts and changes ($O_i, R_i, S_i, I_i, \Delta_i, \rho_i$) and patch volume ($D_i$), together with metadata and diagnostics (Table~\ref{tab:per_file_eval_metrics}).

At the model level we report mean and median discharge, and a workload-weighted discharge:
\[
\rho_w =
\frac{\sum_{i=1}^{N} S_i}{\sum_{i=1}^{N} O_i}\times 100,
\]
which answers: ``Across all baseline obligations, what fraction were discharged?''
We summarize patch closure using WDC and MFDC, and report concentration statistics (median, P95, top-five share) to expose heavy-tailed residual work.

We interpret $I_i$ and $\Delta_i$ only as contract regressions, meaning new triggers under the pinned policy.
They are not behavioral regressions.
Contract metrics are therefore read together with validity, preservation, and loadability diagnostics.

\subsection{Validity and regression diagnostics}
\label{sec:pce_validity_diagnostics}

Contract metrics can be misleading in the presence of omissions, degenerate rewrites, syntax errors, or include-time failures.
Accordingly, we report a fixed diagnostic suite for every output to contextualize contract outcomes.
These checks provide signals of validity and regressions, but they do not establish functional equivalence.

\paragraph{Diagnostic tiers}
We report:
(i) syntax validation via \texttt{php -l} (PHP~8.3),
(ii) structural preservation signals (Table~\ref{tab:structural_completeness_signals}),
(iii) contract metrics under Rector,
(iv) PHPCompatibility results, and
(v) a loadability sentinel (Section~\ref{sec:pce_loadability}).
Outcome definitions appear in Table~\ref{tab:outcome_matrices}.
Metrics that require parsing are N/A when \texttt{php -l} fails.

\subsection{Isolated loadability sentinel}
\label{sec:pce_loadability}

We run a loadability sentinel on a fixed slice ($N_{\mathrm{rt}}=55$).
Each file is loaded once via \texttt{require\_once} in a minimal WordPress stub harness.
Parse errors, fatal errors, and uncaught exceptions count as failures.
We interpret this tier only as a regression signal.
A regression is orig pass and migrated fail.

% =========================
% Results table 
% =========================
\begin{table*}[t]
\centering
\caption[Evaluation summary under pinned contract]{Pinned-contract evaluation on the coverage-optimized diagnostic stress test. All quantities are contract-relative. Panel A reports the primary outcome (incidence-level discharge). Panel B stratifies discharge by file size. Panel C reports discharge by PHP family. Patch-volume closure is reported for localization only (not effort).}
\label{tab:results_panels_revised}
\footnotesize
\setlength{\tabcolsep}{4pt}
\renewcommand{\arraystretch}{1.12}

% ---------------- Panel A ----------------
\textbf{A. Obligation discharge (incidence-level breadth; pinned contract)}\par
\vspace{2pt}
\resizebox{\textwidth}{!}{%
\begin{tabular}{lrrrrrrrr}
\toprule
\textbf{System} &
\textbf{Mean $\rho$ (\%)} &
\textbf{Weighted $\rho_w$ (\%)} &
\textbf{Discharged} &
\textbf{Remaining} &
\textbf{Introduced} &
\textbf{Net $\Delta$} &
\textbf{complete} &
\textbf{$\rho_i<50\%$} \\
\midrule
Rector-applied (reference)   & 98.55 & 98.43 & 503 & 17  & 9 & 494 & 89 & 0  \\
\midrule
Gemini 2.5 Pro               & 77.11 & 75.76 & 372 & 140 & 21 & 351 & 19 & 3  \\
Gemini 2.5 Flash             & 62.32 & 59.48 & 273 & 203 & 17 & 256 & 13 & 24 \\
GPT-5 Codex                  & 71.35 & 69.39 & 297 & 151 & 20 & 277 & 15 & 12 \\
Claude Sonnet 4 (2025-05-14) & 56.97 & 54.51 & 272 & 244 & 17 & 255 & 6  & 33 \\
Llama 3.3 70B Instruct       & 34.21 & 32.16 & 156 & 347 & 18 & 138 & 2  & 66 \\
\bottomrule
\end{tabular}%
}
\vspace{2pt}
\noindent\footnotesize Aggregates are computed over analyzable files per model after excluding migrated files where Rector dry-run failed: Rector 100, Claude 97, Gemini Flash 91, Gemini Pro 97, GPT-5 Codex 88, Llama 95.\normalsize

% ---------------- Panel B ----------------
\textbf{B. Obligation discharge by file size (Mean $\rho$; discharged/total incidences)}\par
\vspace{2pt}
{\scriptsize
\resizebox{\textwidth}{!}{%
\begin{tabular}{lcccc}
\toprule
\textbf{Model} &
\textbf{Small} &
\textbf{Medium} &
\textbf{Large} &
\textbf{Extra-large} \\
\midrule
Gemini 2.5 Pro &
77.63 (77/99) & 83.81 (136/162) & 74.76 (105/141) & 60.58 (54/89) \\
Gemini 2.5 Flash &
71.51 (70/99) & 66.97 (104/153) & 51.22 (67/131)  & 41.84 (32/76) \\
GPT-5 Codex &
75.86 (76/99) & 75.88 (111/146) & 66.58 (86/130) & 45.13 (24/53) \\
Claude Sonnet 4 (2025-05-14) &
65.75 (60/91) & 69.33 (111/158) & 41.53 (63/146)  & 38.30 (38/104) \\
Llama 3.3 70B Instruct &
37.76 (35/93) & 42.70 (71/162) & 28.50 (40/140) & 11.85 (10/90) \\
\bottomrule
\end{tabular}%
}}
\vspace{2pt}

% ---------------- Panel C ----------------
\noindent\textbf{C. Discharged obligations by PHP family (analyzable subset per model; discharged/total; \% discharged)}\par
\vspace{2pt}
\begin{tabular*}{\textwidth}{@{\extracolsep{\fill}}p{0.30\textwidth}ccc@{}}
\toprule
\textbf{Model} & \textbf{PHP 5.x} & \textbf{PHP 7.x} & \textbf{PHP 8.x} \\
\midrule
Gemini 2.5 Flash & 116/137 (84.67\%) & 97/120 (80.83\%) & 60/202 (29.70\%) \\
Gemini 2.5 Pro   & 143/149 (95.97\%) & 114/126 (90.48\%) & 115/216 (53.24\%) \\
GPT-5 Codex      & 126/133 (94.74\%) & 88/106 (83.02\%) & 83/189 (43.92\%) \\
Claude Sonnet 4 (2025-05-14) & 109/150 (72.67\%) & 92/129 (71.32\%) & 71/220 (32.27\%) \\
Llama 3.3 70B Instruct & 87/145 (60.00\%) & 53/125 (42.40\%) & 16/215 (7.44\%) \\
\bottomrule
\end{tabular*}
\vspace{4pt}

% --------- Supplementary ----------
\noindent\hrule height 0.4pt
\vspace{4pt}
\noindent\textbf{Secondary: Patch-volume closure (pinned contract; localization only)}\par
\vspace{2pt}
\begin{tabular*}{\textwidth}{@{\extracolsep{\fill}}p{0.30\textwidth}rrrr@{}}
\toprule
\textbf{System} &
\textbf{WDC (\%)} &
\textbf{MFDC (\%)} &
\textbf{Top-5 share (\%)} &
\textbf{Remaining patch lines} \\
\midrule
Rector-applied (reference)   & 95.08 & 98.19 & 91.43 & 1{,}109  \\
GPT-5 Codex                  & 67.46 & 69.93 & 36.46 & 5{,}074  \\
Gemini 2.5 Pro               & 71.92 & 76.90 & 50.70 & 5{,}826  \\
Gemini 2.5 Flash             & 60.65 & 64.75 & 43.53 & 7{,}620  \\
Claude Sonnet 4 (2025-05-14) & 57.77 & 62.75 & 40.31 & 9{,}411  \\
Llama 3.3 70B Instruct       & 36.71 & 42.33 & 39.35 & 12{,}342 \\
\bottomrule
\end{tabular*}
\end{table*}

\subsection{Cross-instrument triangulation with PHPCompatibility}
\label{sec:pce_phpcompat}

We triangulate using PHPCompatibility.\cite{phpcompat}
For each output set we run PHPCS with \texttt{--standard=PHPCompatibility} and \texttt{--runtime-set testVersion 8.3}.
We parse JSON reports and compute issue reduction:
\[
\text{IssueReduction}(\%)=
\frac{\text{Issues}_{\mathrm{orig}}-\text{Issues}_{\mathrm{model}}}
{\text{Issues}_{\mathrm{orig}}}\times 100.
\]
PHPCompatibility uses a different rule set and granularity than Rector.
Agreement suggests consistent movement under both policies.
Divergence reflects policy differences rather than a contradiction.

\subsection{Pipeline completeness and residual concentration}
\label{sec:pce_completeness_concentration}

We report:
\begin{align}
N_{\mathrm{contract\text{-}clean}} &= |\{ i \mid R_i=0 \}|,\\
N_{\mathrm{low\text{-}compliance}} &= |\{ i \mid \rho_i<50 \}|.
\end{align}
$N_{\mathrm{contract\text{-}clean}}$ counts files with no remaining triggers and captures breadth only.
$N_{\mathrm{low\text{-}compliance}}$ counts files with discharge below $50\%$.
We also report concentration over remaining triggers and patch volume (median, P95, top-five share).

\subsection{Contract-aligned reference baseline}
\label{sec:pce_baseline}

We include a contract-aligned reference baseline, \emph{Rector applied}.
We apply the pinned ruleset to each file, then re-run the same contract on the rewritten output.

This baseline is not guaranteed to be trigger-free.
Rector rules are conservative and pattern-limited.
It provides a practical ceiling under the same policy and scoring protocol.